\subsection{Final Experimental Recommendation}
\label{sec:final-configuration}

The experiments support retaining the original masked, pocket-filtered Hough
detector and recommending the tuned K=6 recognition configuration with the
solid-only CCA correction. This is the strongest observed combination of
localisation, type and number recognition among the evaluated alternatives.
The recommendation is based on pooled results and recorded regressions,
rather than selecting a method because it recovers one visually convincing
example. Table1--Table6 were used to tune recognition, so the recommended
recognition scores remain exploratory.

\begin{table}[htbp]
\centering
\small
\caption{Recognition progression with the original Hough detections fixed.
All counts are over the same 86 matched balls in Table1--Table6.}
\label{tab:final-recognition-progression}
\begin{tabular}{lrrrr}
\toprule
Configuration & Correct type & Type accuracy & Correct number & Number accuracy \\
\midrule
Original K=2 baseline & 76/86 & 88.37\% & 67/86 & 77.91\% \\
Tuned K=6 & 81/86 & 94.19\% & 74/86 & 86.05\% \\
Tuned K=6 + solid-only CCA & 82/86 & 95.35\% & 75/86 & 87.21\% \\
\bottomrule
\end{tabular}
\end{table}
\FloatBarrier

The recommended settings are:
\begin{itemize}
\item \textbf{Detection:} retain $dp=1.5$, minimum centre distance 18,
upper Canny threshold 75, accumulator threshold 35, radii 15--30,
the known table mask and the existing 25-table-pixel pocket exclusion.
\item \textbf{Extraction:} retain the detected radius, with zero padding
and no cloth filtering.
\item \textbf{Clustering:} use six BGR colour clusters with the existing
fixed-seed K-means implementation.
\item \textbf{Light definition:} require $V\geq170$ and $S\leq125$,
consistently in type classification, body-cluster selection and CCA.
\item \textbf{Base type:} require at least 85\% light for the cue ball
or at least 80\% non-light for a solid; otherwise retain the existing
striped/unknown rules.
\item \textbf{Number recognition:} use blue hue $100\leq H<125$
and purple $125\leq H<165$; retain the other original number thresholds
and the minimum eligible cluster area of 2\%.
\item \textbf{CCA:} retain eight-connected components above a 2\%
ball-area cutoff and a $5\times5$ erosion boundary band. Enable only
the solid override: at least 70\% non-light pixels, at most 20\%
light-boundary coverage and exactly one eligible non-light component.
Disable cue and stripe overrides; otherwise retain the base type.
\end{itemize}

Relative to the original K=2 baseline, this combination increased correct
types by six and correct numbers by eight. Of these gains, the targeted
CCA correction accounted for only one additional correct type and number
beyond K=6. On AllBalls, the combination retained the original baseline's
16/16 correct types and numbers. Considering all 90 comparison annotations,
including the four missed detections, correct types corresponded to 82/90
($91.11\%$) and correct numbers to 75/90 ($83.33\%$). These are
detected-and-correct recall measures, not precision measures: false
detections must still be reported separately.

The other experiments informed the following exclusions. Padding admitted
additional background and generally worsened number recognition. Cloth
filtering slightly improved some original-threshold results, but was not
part of the best K=6 combination and could remove real green-ball pixels.
The original full CCA rule set reduced accuracy; only the tuned solid
correction was retained in the recommendation. The distance-peak watershed
prototype recovered no detections. Segmentation-first Hough recovered
three previously missed annotation identities, including ball 10 in
Table3 and Table4, but lost three others and increased false detections
from two to 16. Its watershed check reduced false detections to 11 but
also reduced true positives from 86 to 83. These alternatives therefore
did not justify replacing the standard detector.

The subsequent histogram-backprojection verifier was screened on AllBalls
and evaluated with a frozen conservative setting. It retained every candidate
in Table1--Table6, preserving recall but removing neither false detection.
It was therefore not added to the recommendation. This candidate-verification
result is distinct from cloth-pixel removal during extraction.

There is a distinction between an experimental recommendation and application
integration. The main application's current configuration still requests
K=2 and calls the original type classifier, using the original light and
solid thresholds. It does not invoke the selected solid-only CCA correction.
The K=6 and CCA results above were implemented and measured in the independent
experiment executable. No claim is made that the main application already
produces the recommended recognition scores.

The next practical step is consolidation and evaluation on fresh images,
rather than assuming that another processing stage will necessarily improve
the result. This report records the selected experimental configuration,
its limits, and the unsuccessful alternatives so that later integration
can be checked against the same settings and evidence.
